\section{Computational certificates and formal verification}\label{sec:verification}
The positive coloring certificates are checked using integer arithmetic and
the coordinate definition of orthogonality. The full checker imports neither
the search program nor a SAT solver. It establishes enumeration completeness
from distinctness and Gaussian counts, then checks every pair of vertices.
Three deliberately invalid certificates, respectively omitting a vertex,
duplicating a subspace, and identifying the colors of adjacent coordinate
lines, are rejected. The line-coloring checker independently verifies its
partition and all within-class pairs. For the twelve-color certificate,
\path{develop/check_line_chromatic_exact.py} checks all $1953$ distinct
vertex pairs, including all $961$ edges. It independently reproduces every
clause of the normalized search encoding and its SHA-256 digest, verifies
the assignment against those clauses, and rejects duplicate-vertex,
zero-vector and monochromatic-edge controls.

The radical/quotient clique theorem has an existing Lean formalization.
Its source is \path{formal/NikandishClique.lean}. The theorem is
\texttt{NikandishClique.result}, inside the namespace
\begin{center}
\texttt{D5.S3.Combinatorics.Orthogonality}.
\end{center}
The upstream revision and source hash are recorded in
\path{formal/UPSTREAM.json}; dependencies must be obtained from that pinned
upstream checkout.

The six-dimensional development uses vectors in $\operatorname{Fin}(64)$,
XOR addition, and the coordinate dot product modulo two. A vertex is a
zero-containing, XOR-closed subset other than $\{0\}$, exactly the nonzero
binary subspaces in this representation. The formal graph is simple and
uses coordinate orthogonality between distinct such subsets. The catalogue
contains $2\,825$ entries including zero and $180\,800$ adjoining-vector
witnesses. The theorem \texttt{Catalogue.complete} derives coverage from
the zero entry and closure under adjoining any vector; it does not assume
the Gaussian enumeration counts. The theorem \texttt{Catalogue.colorable}
transfers the checked catalogue coloring to this graph. A checked
$15$-clique supplies the converse inequality in
\texttt{BinaryOrthogonality.chromatic\_number\_six}.

The archived Lean 4.33.0 development verifies four finite assertions with
\texttt{native\_decide}. The final theorem depends on the standard axioms
\texttt{propext}, \texttt{Classical.choice}, and \texttt{Quot.sound}, together
with four native-evaluation axioms; its computational trust therefore includes
native evaluation. The generic completeness and colorability lemmas depend
only on the three standard axioms. The recovered original log corrects the
incomplete historical summary; see \path{notes/formal-verification-audit.md}.
The exhaustive integer certificate independently verifies the same coloring.
The weighted line bound, twelve-color line certificate, geometric profile
counts and dimension-seven reduction
are not covered by this Lean theorem.

The public package~\cite{Package} contains the literal witnesses, independent
checkers, search records, and SHA-256 bindings. The new geometric profiler
checks every assigned color against its $H(U)$-determined list. The
dimension-seven reduction is also replayed in dimension six, where it
reproduces the established core size and Boolean encoding counts.

The dimension-seven lower bound in Section~\ref{sec:seven} has a
written proof. The checker
\path{develop/check_dimension_seven_nineteen_obstruction.py}
independently verifies its coordinate ingredients: all $4096$ quadratic
polarization pairs, the $36/28$ quadratic value counts, all $288$ even
heptads and $2016$ six-point cliques with their unique heptad completions,
and all $8505$ point--hyperplane intersections with the $135$ Lagrangians.
Nine explicit spread controls check the hole-clique obstruction.
These controls do not enumerate all eight-member partial spreads;
Lemma~\ref{lem:eight-completion} proves their completion uniformly.

The additional nineteen-color constraints in Section~\ref{sec:higher}
have their own evidence. The six- and seven-member partial-spread
completion lemmas use written normalization and exhaustive finite
enumerations. The four marked-form exclusions use independently checked
covering certificates. For the zero-five-even-sum form,
\path{develop/check_dimension_seven_five_four_Z_cover.py}
reconstructs the complete necessary-root family and orbit partition.
It then verifies $14\,337$ reachable failed states, all $3517$ branches
and $10\,995$ leaves, using all $224$ anchored heptads. Every branch
removes seven points, leaves admit no row, and the empty set is never
declared failed; induction proves the exclusion. Three invalid
certificates are rejected. The proof notes, checkers and hash-bound reports
are recorded separately in the public package. The small accounting
auditor \path{develop/check_dimension_seven_profile_scope.py}
reconstructs the parity-count catalogs and checks the scope of the
earlier count candidates, including the $64$ rows to which the common
hole capacity applied. Its archived scope predates the completed
four-form exclusion. The Z covering auditor independently reconstructs
the raw count catalogs and checks the removal of exactly the row
$(5,4;6,2,8)$ from the earlier remaining list. The separate
earlier covering results are hash-bound inputs to this combined profile
conclusion; their proofs are not rerun or used in the Z covering proof.
For the further $(3,6;6,2,8)$ exclusion,
\path{develop/check_dimension_seven_three_six_cover.py}
reconstructs all three complete root families and their marked
orbit partitions, then checks all $1279$ representatives using
the full $224$-row anchored-heptad catalog. Its new certificate
has $1949$ reachable failed states, $676$ branches and $1332$
leaves. All pivot successors are checked, every branch removes
seven points, and three invalid controls are rejected. No earlier
covering certificate is a premise of this new finite proof.
The same auditor reconstructs the raw count catalogs and removes
exactly $(3,6;6,2,8)$ from the preceding $26$-row list. There are
then $25$ five-plus-six and $39$ three-six candidates: $62$ have
six through eight pure odd heptads and two have five.
For $(2,6,6;7,0,8)$,
\path{develop/check_dimension_seven_two_six_six_cover.py}
reconstructs all $384$ necessary sets and their $75$ marked orbits,
retaining both actual defect holes and intrinsic pure-six roles.
It checks the full $288$-row even-heptad catalog, including rows
avoiding $M$. The new certificate has $368$ reachable failed states,
$293$ branches and $225$ leaves; all pivot successors are checked,
and four invalid controls are rejected. Induction on the remaining
cardinality proves the outside-$M$ obstruction. The inside-$M$
recoloring exclusion invokes the earlier $(2,6;7,1,8)$ theorem and
its historical finite cover input, recorded in
\path{notes/dimension-seven-five-six-cover.md}; that certificate is
not a premise of the new outside-$M$ proof and is not rerun here.
The new auditor removes exactly this three-six row, preserving
the five-plus-six list. The current totals are $25$ and $38$:
$61$ have six through eight pure odd heptads and two have five.
None of these dimension-seven results is covered by the archived
six-dimensional Lean theorem.
